%% Generated by Sphinx.
\def\sphinxdocclass{report}
\documentclass[a4paper,11pt,oneside,english,openany]{sphinxmanual}
\ifdefined\pdfpxdimen
   \let\sphinxpxdimen\pdfpxdimen\else\newdimen\sphinxpxdimen
\fi \sphinxpxdimen=.75bp\relax
\ifdefined\pdfimageresolution
    \pdfimageresolution= \numexpr \dimexpr1in\relax/\sphinxpxdimen\relax
\fi
%% let collapsible pdf bookmarks panel have high depth per default
\PassOptionsToPackage{bookmarksdepth=5}{hyperref}

\PassOptionsToPackage{booktabs}{sphinx}
\PassOptionsToPackage{colorrows}{sphinx}

\PassOptionsToPackage{warn}{textcomp}
\usepackage[utf8]{inputenc}
\ifdefined\DeclareUnicodeCharacter
% support both utf8 and utf8x syntaxes
  \ifdefined\DeclareUnicodeCharacterAsOptional
    \def\sphinxDUC#1{\DeclareUnicodeCharacter{"#1}}
  \else
    \let\sphinxDUC\DeclareUnicodeCharacter
  \fi
  \sphinxDUC{00A0}{\nobreakspace}
  \sphinxDUC{2500}{\sphinxunichar{2500}}
  \sphinxDUC{2502}{\sphinxunichar{2502}}
  \sphinxDUC{2514}{\sphinxunichar{2514}}
  \sphinxDUC{251C}{\sphinxunichar{251C}}
  \sphinxDUC{2572}{\textbackslash}
\fi
\usepackage{cmap}
\usepackage[T1]{fontenc}
\usepackage{amsmath,amssymb,amstext}
\usepackage{babel}





\usepackage[Bjarne]{fncychap}
\usepackage{sphinx}

\fvset{fontsize=auto}
\usepackage{geometry}


% Include hyperref last.
\usepackage{hyperref}
% Fix anchor placement for figures with captions.
\usepackage{hypcap}% it must be loaded after hyperref.
% Set up styles of URL: it should be placed after hyperref.
\urlstyle{same}

\addto\captionsenglish{\renewcommand{\contentsname}{Contents:}}

\usepackage{sphinxmessages}
\setcounter{tocdepth}{1}


		\renewcommand{\familydefault}{\rmdefault}
        \usepackage{setspace} % Required for more spacing options
		\parindent 0px
		\setstretch{1.1}
        \usepackage[T1]{fontenc}
        \usepackage{charter}
        \usepackage{amsmath}
        \usepackage{amsfonts}
        \usepackage{amssymb}
        \usepackage{fancyhdr}
        \usepackage{arydshln}
        \usepackage{xcolor}
        \usepackage{makeidx}  % To create an index
        \usepackage{tocloft}  % To customize Table of Contents

        % Fancy header and footer
        \pagestyle{fancy}
        \fancyhead[L]{\leftmark}
        \fancyhead[R]{\thepage}
        \fancyfoot[C]{}

        % Table of Contents: Formatting (removing subsections)
        \setcounter{tocdepth}{1}  % Limit ToC to sections only (remove subsections)
        \renewcommand{\cftsecfont}{\normalfont\bfseries}  % Bold section titles
        \renewcommand{\cftsecindent}{0em}                 % No indentation for sections
        \setlength{\cftbeforesecskip}{0.5ex}              % Space before section titles

        % Index Configuration (simplified to sections only)
        \makeindex  % Generate index

        % Page layout adjustments
		\usepackage{geometry}
        \geometry{top=1.5cm, bottom=1.5cm, left=1.5cm, right=1.5cm}
        \usepackage{hyperref}
        \hypersetup{colorlinks=true,linkcolor=violet,filecolor=magenta,urlcolor=blue}

        % Better verbatim environment for code blocks
        %\usepackage{fvextra}
        \usepackage{tcolorbox}
        
        \newtcbox{\tbox}{colframe=black,colback=gray!5,boxrule=0.7pt,arc=4pt,boxsep=5pt}
    

\title{SIMVIS}
\date{May 20, 2026}
\release{1.0.0}
\author{SC,AE,SB,BB}
\newcommand{\sphinxlogo}{\vbox{}}
\renewcommand{\releasename}{Release}
\makeindex
\begin{document}

\ifdefined\shorthandoff
  \ifnum\catcode`\=\string=\active\shorthandoff{=}\fi
  \ifnum\catcode`\"=\active\shorthandoff{"}\fi
\fi

\pagestyle{empty}
\sphinxmaketitle
\pagestyle{plain}
\sphinxtableofcontents
\pagestyle{normal}
\phantomsection\label{\detokenize{index::doc}}
\sphinxstepscope



\sphinxstepscope


\chapter{simvis module}
\label{\detokenize{simvis:module-simvis}}\label{\detokenize{simvis:simvis-module}}\label{\detokenize{simvis::doc}}\index{module@\spxentry{module}!simvis@\spxentry{simvis}}\index{simvis@\spxentry{simvis}!module@\spxentry{module}}

\section{All sky simulations of Gaussian random sky signal and corresponding Visibilities}
\label{\detokenize{simvis:all-sky-simulations-of-gaussian-random-sky-signal-and-corresponding-visibilities}}
\sphinxAtStartPar
\sphinxstylestrong{For a given Angular Power Spectrum {[}APS, in} \(\text{mK}^2\) \sphinxstylestrong{unit {]} of brightness temperature fluctuations or Power Spectrum {[}PS, in}   \(\text{mK}^2\;\text{Mpc}^3\) \sphinxstylestrong{unit {]} and an input UVFITS file, this outputs a new UVFITS file with simulated visibilities.}


\bigskip\hrule\bigskip


\sphinxAtStartPar
\sphinxstylestrong{Our aim is to simulate the visibilities for a given telescope and for either a given Angular Power Spectrum  or a given Power Spectrum.}

\sphinxAtStartPar
For given Angular Power Spectrum \(C_{\ell}\), we generate muliple realizations of the sky signal at the centered frequency \(\nu_c\), and for a given Power Spectrum \(P(k)\), we generate 3D sky signal for all frequency channels (frequency channel information taken from input \sphinxcode{\sphinxupquote{UVFITS}} file). Afterwards by this simulated sky signal, we simulate the corresponding visibilities.

\sphinxAtStartPar
Visibilities that are simulated are in Jy(Jansky) unit. The visibility at a baseline \(\vec{\textbf{U}}\) is given by:
\begin{equation*}
\begin{split}\mathcal{V}(\vec{\textbf{U}},\nu)=Q_\nu \int_{\rm UH}\ d\Omega_{\hat{n}}\ T(\hat{\mathbf{n}},\nu)\,{A}(\Delta\hat{\mathbf{n}},\nu)\ \textbf{exp}\left[{2\pi \text{i}\vec{\textbf{U}}\cdot\Delta\hat{\mathbf{n}}}\right]\end{split}
\end{equation*}\begin{itemize}
\item {} 
\sphinxAtStartPar
\(T(\hat{\mathbf{n}},\nu)\) is the surface brightness temperature field.

\item {} 
\sphinxAtStartPar
\({A}(\Delta\hat{\mathbf{n}},\nu)\) is the primary beam pattern of the telescope.

\item {} 
\sphinxAtStartPar
\(\textbf{exp}\left[{2\pi \text{i}\vec{\textbf{U}}\cdot\Delta\hat{\mathbf{n}}}\right]\) is the phase factor.

\end{itemize}

\sphinxAtStartPar
In \sphinxcode{\sphinxupquote{Healpix}} we discretize the sky and the integral becomes summation. But the integral is restricted only to  upper hemisphere.

\sphinxAtStartPar
\sphinxstylestrong{The discritized version is given as}:
\begin{equation*}
\begin{split}\mathcal{V}(\alpha_{p},\vec{\textbf{U}},\nu) = Q_{\nu}\ \Delta\Omega_{pix}\ \sum_{q} \ T(\hat{\mathbf{n}}_{q},\nu) \ {A}(\Delta\hat{\mathbf{n}}_{q},\nu) \ \textbf{exp}\left[{2\pi \text{i}\vec{\textbf{U}}\cdot\Delta\hat{\mathbf{n}}_{q}}\right]\end{split}
\end{equation*}
\sphinxAtStartPar
where \(\alpha_p\) is the pointing RA. For more details, \sphinxhref{https://arxiv.org/pdf/2212.01251}{see Section (2)}.

\sphinxAtStartPar
Also,
\begin{align*}\!\begin{aligned}
Q_{\nu}  &=2k_B /\lambda^2\\
\Delta\hat{\mathbf{n}}&= \hat{\mathbf{n}}- \hat{\mathbf{p}}\\
\vec{\textbf{U}} &= u\;{\hat{\mathbf{e}}_ 1 (\alpha_{p})} + v \;{\hat{\mathbf{e}}_ 2 (\alpha_{p})} + w \;{\hat{\mathbf{e}}_3(\alpha_{p})}\\
\end{aligned}\end{align*}
\sphinxAtStartPar
\(\hat{\mathbf{e}}_1,\hat{\mathbf{e}}_2,\hat{\mathbf{e}}_3\) are the \sphinxstylestrong{basis} vectors defined in the local tangent plane. \(\hat{\mathbf{e}}_1\) is along \sphinxstylestrong{east}, \(\hat{\mathbf{e}}_2\) is along \sphinxstylestrong{north} and \(\hat{\mathbf{e}}_3\) is \sphinxstylestrong{vertically overhead, i.e. zenith pointed}.

\sphinxAtStartPar
The \sphinxstylestrong{basis} vectors \(\hat{\mathbf{e}}_1,\hat{\mathbf{e}}_2,\hat{\mathbf{e}}_3\) are dependent on pointing \sphinxstylestrong{RA} \(\alpha_p\) and \sphinxstylestrong{declination(DEC)} \(\delta_p\). Components of basis vectors in the cartesian system \(xyz\) is given by:
\begin{equation*}
\begin{split}\hat{\mathbf{e}}_1&=[-\sin{\alpha_p},\cos{\alpha_p},0]\\
\hat{\mathbf{e}}_2&=[-\cos{\alpha_p}\sin{\delta_p},-\sin{\alpha_p}\sin{\delta_p},\cos{\delta_p}]\\
\hat{\mathbf{e}}_3&=[\cos{\alpha_p}\cos{\delta_p},\sin{\alpha_p}\cos{\delta_p},\sin{\delta_p}]\end{split}
\end{equation*}
\sphinxAtStartPar
Where \(\Delta\Omega_{pix}\) refers to the solid angle subtended by each simulation pixel and \(\hat{\mathbf{p}}\) refers to the pointing direction of the antenna (here zenith pointed), and hence \(\hat{\mathbf{p}}=\hat{\mathbf{e}}_3\).
\subsubsection*{Notes}

\sphinxAtStartPar
In the simulation of visibilities for 3D sky signal, we have not incorporated the fact that the baselines \(\vec{\textbf{U}}\) and antenna beam pattern \({A}(\Delta\hat{\mathbf{n}},\nu)\) changes as we vary the frequency. We assume only the brightness temperature field \(T(\hat{\mathbf{n}},\nu)\) changes with frequency.


\section{To run:}
\label{\detokenize{simvis:to-run}}
\begin{sphinxVerbatim}[commandchars=\\\{\}]
\PYG{o}{\PYGZgt{}\PYGZgt{}} \PYG{k+kn}{import} \PYG{n+nn}{simvis} \PYG{k}{as} \PYG{n+nn}{sv}
\PYG{o}{\PYGZgt{}\PYGZgt{}} \PYG{c+c1}{\PYGZsh{} skysimtype = \PYGZsq{}2D\PYGZsq{} for 2D sky signal (multiple realizations)}
\PYG{o}{\PYGZgt{}\PYGZgt{}} \PYG{c+c1}{\PYGZsh{} skysimtype = \PYGZsq{}3D\PYGZsq{} for 3D sky signal}
\PYG{o}{\PYGZgt{}\PYGZgt{}} \PYG{n}{sv}\PYG{o}{.}\PYG{n}{sim\PYGZus{}vis}\PYG{p}{(}\PYG{n}{infits}\PYG{p}{,} \PYG{n}{outfits}\PYG{p}{,} \PYG{n}{skysimtype}\PYG{p}{)}
\PYG{o}{\PYGZgt{}\PYGZgt{}} \PYG{c+c1}{\PYGZsh{} If you want the seed to be, say seed = 5, run:}
\PYG{o}{\PYGZgt{}\PYGZgt{}} \PYG{n}{sv}\PYG{o}{.}\PYG{n}{sim\PYGZus{}vis}\PYG{p}{(}\PYG{n}{infits}\PYG{p}{,} \PYG{n}{outfits}\PYG{p}{,} \PYG{n}{skysimtype}\PYG{p}{,} \PYG{n}{seed} \PYG{o}{=} \PYG{l+m+mi}{5}\PYG{p}{)}  \PYG{c+c1}{\PYGZsh{} seed = None for random.}
\PYG{o}{\PYGZgt{}\PYGZgt{}} \PYG{c+c1}{\PYGZsh{} for 2D sky signal}
\PYG{o}{\PYGZgt{}\PYGZgt{}} \PYG{n}{sv}\PYG{o}{.}\PYG{n}{sim\PYGZus{}vis}\PYG{p}{(}\PYG{n}{infits}\PYG{p}{,} \PYG{n}{outfits}\PYG{p}{,} \PYG{n}{skysimtype} \PYG{o}{=} \PYG{l+s+s1}{\PYGZsq{}}\PYG{l+s+s1}{2D}\PYG{l+s+s1}{\PYGZsq{}}\PYG{p}{,} \PYG{n}{seed}\PYG{p}{)}
\PYG{o}{\PYGZgt{}\PYGZgt{}} \PYG{c+c1}{\PYGZsh{} for 3D sky signal}
\PYG{o}{\PYGZgt{}\PYGZgt{}} \PYG{n}{sv}\PYG{o}{.}\PYG{n}{sim\PYGZus{}vis}\PYG{p}{(}\PYG{n}{infits}\PYG{p}{,} \PYG{n}{outfits}\PYG{p}{,} \PYG{n}{skysimtype} \PYG{o}{=} \PYG{l+s+s1}{\PYGZsq{}}\PYG{l+s+s1}{3D}\PYG{l+s+s1}{\PYGZsq{}}\PYG{p}{,}\PYG{n}{seed}\PYG{p}{)}
\end{sphinxVerbatim}


\bigskip\hrule\bigskip


\sphinxAtStartPar
\sphinxstylestrong{In generating 2D sky signal (for multiple realizations)}, for each realization of GRF, a seed is given. By default seed is set to \sphinxstylestrong{None}.

\sphinxAtStartPar
If you are doing \(N_{\rm realizations}\), in each realization seed is incremented by 1 w.r.t. previous value. For example, say you want to generate 10 sky maps, i.e.  \(N_{\rm realizations}=10\), and you gave seed = 5(say). Then while doing \(1^{\text{st}}\) realization seed will be 5, in next realization seed will increment by 1 i.e. 6 and so on.


\bigskip\hrule\bigskip


\sphinxAtStartPar
\sphinxstylestrong{For 2D sky signal case:}

\sphinxAtStartPar
The entire code works at a single frequency, the centered frequency \(\nu_c\), which it reads from the input UVFITS file. It then simulates the visibilities at a RA,DEC which is read from input UVFITS file.
\begin{itemize}
\item {} 
\sphinxAtStartPar
Mutiple realizations (number of realizations \(N_{\rm realizations}\)) are written along the frequency axis of output FITS.

\item {} 
\sphinxAtStartPar
\(N_{\rm realizations}\) cannot exceed \(N_{\rm channels}\) the number of channels in the original UVFITS.

\item {} 
\sphinxAtStartPar
Channels \(N_{\rm realizations}\) +1 to \(N_{\rm channels}\) of are left untouched in the output FITS file.

\item {} 
\sphinxAtStartPar
The maximum value of angular multiplole \(\ell_{max}\) is by default \((3\times nside)-1\).

\end{itemize}

\newpage

\sphinxAtStartPar
\sphinxstylestrong{For 3D sky signal case}:
\begin{itemize}
\item {} 
\sphinxAtStartPar
Only the sky signal is frequency dependent.

\item {} 
\sphinxAtStartPar
The sky signal and visibilities are being simulated for all the channels present in the input \sphinxcode{\sphinxupquote{UVFITS}} file.

\item {} 
\sphinxAtStartPar
Frequency dependent antenna beam patten and baselines are not incorporated.

\end{itemize}


\bigskip\hrule\bigskip


\sphinxAtStartPar
If you want to do some operations with your data in uvinfits file but do not wish to write it in a uvoutfits file you can call the function \sphinxcode{\sphinxupquote{sim\_vis\_hdul}}. And it will update the data in the hdulist.

\sphinxAtStartPar
\sphinxcode{\sphinxupquote{For example run :}}

\begin{sphinxVerbatim}[commandchars=\\\{\}]
\PYG{o}{\PYGZgt{}\PYGZgt{}} \PYG{k}{with} \PYG{n}{fits}\PYG{o}{.}\PYG{n}{open}\PYG{p}{(}\PYG{n}{uvinfits}\PYG{p}{)} \PYG{k}{as} \PYG{n}{hdulist}\PYG{p}{:}
\PYG{o}{\PYGZgt{}\PYGZgt{}}      \PYG{n}{sv}\PYG{o}{.}\PYG{n}{sim\PYGZus{}vis\PYGZus{}hdul}\PYG{p}{(}\PYG{n}{hdulist}\PYG{p}{,} \PYG{n}{skysimtype}\PYG{p}{,} \PYG{n}{seed} \PYG{p}{)}
\PYG{o}{\PYGZgt{}\PYGZgt{}} \PYG{c+c1}{\PYGZsh{} put your operations or code function here that will act on hdulist.}
\PYG{o}{\PYGZgt{}\PYGZgt{}} \PYG{c+c1}{\PYGZsh{} if you want to save the data use :}
\PYG{o}{\PYGZgt{}\PYGZgt{}} \PYG{n}{hdulist}\PYG{o}{.}\PYG{n}{writeto}\PYG{p}{(}\PYG{n}{uvoutfits}\PYG{p}{,}\PYG{n}{overwrite} \PYG{o}{=} \PYG{k+kc}{True}\PYG{p}{)}
\end{sphinxVerbatim}


\bigskip\hrule\bigskip



\subsection{\sphinxstylestrong{Angular Power Spectrum Function}}
\label{\detokenize{simvis:angular-power-spectrum-function}}
\sphinxAtStartPar
If you want to give your own APS function as in the input please make the APS function first and make sure it has only \(\ell\) dependence. Say your APS function is myaps(l) and you want to generate the GRF by this APS.
To do this run:
\begin{quote}

\begin{sphinxVerbatim}[commandchars=\\\{\}]
\PYG{o}{\PYGZgt{}\PYGZgt{}} \PYG{k+kn}{import} \PYG{n+nn}{numpy} \PYG{k}{as} \PYG{n+nn}{np}
\PYG{o}{\PYGZgt{}\PYGZgt{}} \PYG{k+kn}{import} \PYG{n+nn}{simvis} \PYG{k}{as} \PYG{n+nn}{sv}
\PYG{o}{\PYGZgt{}\PYGZgt{}} \PYG{k}{def} \PYG{n+nf}{myaps}\PYG{p}{(}\PYG{n}{l}\PYG{p}{)}\PYG{p}{:}
\PYG{o}{\PYGZgt{}\PYGZgt{}}     \PYG{n}{Amp}  \PYG{o}{=} \PYG{l+m+mi}{50}   \PYG{c+c1}{\PYGZsh{} in mK\PYGZca{}2 unit}
\PYG{o}{\PYGZgt{}\PYGZgt{}}     \PYG{n}{beta} \PYG{o}{=} \PYG{o}{\PYGZhy{}}\PYG{l+m+mi}{1}   \PYG{c+c1}{\PYGZsh{} power index}
\PYG{o}{\PYGZgt{}\PYGZgt{}}     \PYG{k}{return} \PYG{n}{Amp} \PYG{o}{*} \PYG{n}{np}\PYG{o}{.}\PYG{n}{power}\PYG{p}{(}\PYG{n}{l}\PYG{p}{,} \PYG{n}{beta}\PYG{p}{)}
\PYG{o}{\PYGZgt{}}
\PYG{o}{\PYGZgt{}\PYGZgt{}} \PYG{n}{sv}\PYG{o}{.}\PYG{n}{sim\PYGZus{}vis}\PYG{p}{(}\PYG{n}{infits}\PYG{p}{,} \PYG{n}{outfits}\PYG{p}{,} \PYG{n}{skysimtype} \PYG{o}{=} \PYG{l+s+s1}{\PYGZsq{}}\PYG{l+s+s1}{2D}\PYG{l+s+s1}{\PYGZsq{}} \PYG{p}{,}
                        \PYG{n}{seed} \PYG{o}{=} \PYG{l+m+mi}{5}\PYG{p}{,} \PYG{n}{apsfunc} \PYG{o}{=} \PYG{n}{myaps} \PYG{p}{)}
\end{sphinxVerbatim}
\end{quote}

\sphinxAtStartPar
The pre\sphinxhyphen{}defined APS function \sphinxcode{\sphinxupquote{aps}} has the form of
\begin{equation*}
\begin{split}C_{\ell}=A\, \left(\frac{\ell}{\ell_0}\right)^{\beta}\end{split}
\end{equation*}
\sphinxAtStartPar
A defines the amplitude of the APS is in the unit of \(\text{mK}^2\). \(\beta\) here is the power\sphinxhyphen{}law index. \(C_{\ell}\) has the unit of  (\(\text{mK}^2\)). By default set to \(A = 100\;\text{mK}^2\), \(\ell_0=1\), \(\beta\) = \sphinxhyphen{}2. (See more in module functions).


\bigskip\hrule\bigskip



\subsection{\sphinxstylestrong{Power Spectrum Function}}
\label{\detokenize{simvis:power-spectrum-function}}
\sphinxAtStartPar
Power spectrum function \sphinxcode{\sphinxupquote{ps}} here is defined as :
\begin{equation*}
\begin{split}\begin{align}
P^{M}(k) = A \left(\dfrac{k}{k_0}\right)^{s}
\end{align}\end{split}
\end{equation*}
\sphinxAtStartPar
with \(A=1, k_0=1, s=-2\) (Default set). \(A\) has the unit of \(\text{mK}^2\;\text{Mpc}^3\). \(k, k_0\) are in \(\text{Mpc}^{-1}\) unit.

\sphinxAtStartPar
We generate 3D sky signal (temperature field) across all the given frequency channels by :
\begin{equation*}
\begin{split}\begin{align}
\boxed{T(\hat{\mathbf{n}}, \nu) = \sum_{\ell, m} a_{\ell m}(\nu) \, Y_{\ell m}(\hat{\mathbf{n}})}
\end{align}\end{split}
\end{equation*}
\sphinxAtStartPar
where we generate \(a_{\ell m}\) coefficients at \(n^{\text{th}}\) frequency channel by :
\begin{equation*}
\begin{split}\begin{align}
a_{\ell m}(\nu_n) = \sum_{q=0}^{N_c- 1} \left[
\sqrt{\dfrac{P^M(k_\perp, k_{\parallel q})}{2N_c\;\Delta{\nu_c }\;r^2\;r'}}
\left({\hat{\mathbf{x}}_q + i \hat{\mathbf{y}}_q} \right)
\right] {\textbf{exp}}\left[{\dfrac{2\pi i n q}{N_c}}\right]
\end{align}\end{split}
\end{equation*}
\sphinxAtStartPar
\sphinxstylestrong{Where} \sphinxhyphen{}
\begin{enumerate}
\sphinxsetlistlabels{\arabic}{enumi}{enumii}{}{.}%
\item {} 
\sphinxAtStartPar
\(k_{\perp} = \ell/r\) and \(k_{\parallel}\) are the components of comoving wave vector \(\vec{k}\) which is perpendicular and parallel (along) to the line of sight respectively.

\item {} 
\sphinxAtStartPar
\(r\) is the comoving distance at centered frequency \(\nu_c\).

\item {} 
\sphinxAtStartPar
\(r' = \dfrac{dr}{d\nu}\) evaluated at \(\nu_c\).

\item {} 
\sphinxAtStartPar
\(N_c\) is the number of frequency channels.

\item {} 
\sphinxAtStartPar
\(\Delta{\nu_c }\) is the frequency separation between two consecutive frequency channels.

\item {} 
\sphinxAtStartPar
\(\nu_n\) refers to n\sphinxhyphen{}th frequency channel. So by convention \(\nu_0, \nu_{N_c-1}\) will be the frequencies of first and last frequency channel.

\item {} 
\sphinxAtStartPar
\(\hat{\mathbf{x}}_q,\hat{\mathbf{y}}_q \sim \mathcal{N}(0,1)\) are independent Gaussian random variables of unit variance  and by equating \(a_{\ell m}(\nu_n)\)’s, we generate \(T(\hat{\mathbf{n}}, \nu_n)\) by Eq. (1.2).

\end{enumerate}
\subsubsection*{Notes}

\sphinxAtStartPar
As temperature field is a real field, we must remember that this follows a property:
\begin{equation*}
\begin{split}\boxed{a_{\ell, m} = (-1)^{m} [a_{\ell, -m}]^{*}}\end{split}
\end{equation*}
\sphinxAtStartPar
\(\text{where} * \text{means complex conjugation.}\) From this proprety we have if \(m=0\) then we have \(a_{\ell, m=0} = [a_{\ell, m=0}]^{*}\), meaning \(a_{\ell, m=0}\) must be real.

\sphinxAtStartPar
So for \(m=0\) case we take the real part of the \(a_{\ell, m}\) generated by Eq. (1.3), i.e.
\begin{equation*}
\begin{split}\begin{align}
a_{\ell}^{m=0}(\nu_n) ={\mathcal Re}\left\{ \sum_{q=0}^{N_c- 1} \left[
\sqrt{\dfrac{P^M(k_\perp, k_{\parallel q})}{2N_c\;\Delta{\nu_c }\;r^2\;r'}}
\left({\hat{\mathbf{x}}_q + i \hat{\mathbf{y}}_q} \right)
\right] {\textbf{exp}}\left[{\dfrac{2\pi i n q}{N_c}}\right]\right\}
\end{align}\end{split}
\end{equation*}

\bigskip\hrule\bigskip


\sphinxAtStartPar
We have the relation between co\sphinxhyphen{}moving distance (\(r\)) to the red\sphinxhyphen{}shift (\(z\)) given by:
\begin{quote}
\begin{equation*}
\begin{split}r_{\nu}=\int_{0}^{z} \dfrac{c\;dz'}{H(z')} \implies \dfrac{\partial r_{\nu}}{\partial z} = \dfrac{c}{H(z)}\end{split}
\end{equation*}\end{quote}

\sphinxAtStartPar
Where \(H(z)\) is the Hubble parameter at red\sphinxhyphen{}shift \(z\) given by (In Flat \(\Lambda\text{CDM}\) model):
\begin{equation*}
\begin{split}H(z)=H_0\left[\Omega_{m0}\;(1+z)^3+\Omega_{\Lambda0}\right]^{1/2}\end{split}
\end{equation*}
\sphinxAtStartPar
We define :
\begin{equation*}
\begin{split}r'_{\nu}&=\dfrac{\partial r_{\nu}}{\partial \nu}\\
        &= \dfrac{\partial r_{\nu}}{\partial z}\cdot\dfrac{\partial z}{\partial \nu}\\
        & = -\dfrac{\nu_0}{\nu^2}\cdot \dfrac{c}{H(z)}\\
        & = -\dfrac{(1+z)^2}{\nu_0} \cdot \dfrac{c}{H(z)}\end{split}
\end{equation*}
\sphinxAtStartPar
Where we have used \(\nu=\nu_0/(1+z)\). Implying \(\dfrac{\partial z}{\partial \nu} = -\dfrac{\nu_0}{\nu^2}=-\dfrac{(1+z)^2}{\nu_0}\)

\sphinxAtStartPar
\(\nu_0 = 1420\; \mathrm{MHz}\). We take only the absolute value \(|r'_{\nu}|=\dfrac{(1+z)^2}{\nu_0} \cdot \dfrac{c}{H(z)}\).

\sphinxAtStartPar
\sphinxstylestrong{The cosmological parameters are taken from} \sphinxhref{https://docs.astropy.org/en/latest/api/astropy.cosmology.realizations.Planck18.html}{Planck 18 Astropy} , \sphinxhref{https://docs.astropy.org/en/latest/api/astropy.cosmology.FlatLambdaCDM.html}{FlatLambdaCDM}.
\subsubsection*{Notes}

\sphinxAtStartPar
\sphinxstylestrong{To generate 3D sky signal}, needed parameters \(\textbf{nuc, nc, dnu}\), rp is \(r'\), nc is the no of channels \(N_c\), dnu is channel width \(\Delta\nu_c\), nuc is the centered frequnency \(\nu_c\).

\sphinxAtStartPar
If you want to give your own power spectrum function (make sure it is vectorized)
\begin{quote}

\begin{sphinxVerbatim}[commandchars=\\\{\}]
\PYG{o}{\PYGZgt{}\PYGZgt{}} \PYG{k+kn}{import} \PYG{n+nn}{numpy} \PYG{k}{as} \PYG{n+nn}{np}
\PYG{o}{\PYGZgt{}\PYGZgt{}} \PYG{k+kn}{import} \PYG{n+nn}{builtins} \PYG{k}{as} \PYG{n+nn}{blt}
\PYG{o}{\PYGZgt{}\PYGZgt{}} \PYG{k+kn}{import} \PYG{n+nn}{simvis} \PYG{k}{as} \PYG{n+nn}{sv}
\PYG{o}{\PYGZgt{}\PYGZgt{}} \PYG{c+c1}{\PYGZsh{} make sure all are float.}
\PYG{o}{\PYGZgt{}\PYGZgt{}} \PYG{n}{blt}\PYG{o}{.}\PYG{n}{A}  \PYG{o}{=}  \PYG{l+m+mf}{1.0}
\PYG{o}{\PYGZgt{}\PYGZgt{}} \PYG{n}{blt}\PYG{o}{.}\PYG{n}{s}  \PYG{o}{=} \PYG{o}{\PYGZhy{}}\PYG{l+m+mf}{3.0}
\PYG{o}{\PYGZgt{}\PYGZgt{}} \PYG{n}{blt}\PYG{o}{.}\PYG{n}{k0} \PYG{o}{=}  \PYG{l+m+mf}{1.0}
\PYG{o}{\PYGZgt{}\PYGZgt{}} \PYG{n+nd}{@np}\PYG{o}{.}\PYG{n}{vectorize}
\PYG{o}{\PYGZgt{}\PYGZgt{}} \PYG{k}{def} \PYG{n+nf}{mypkfunc}\PYG{p}{(}\PYG{n}{k}\PYG{p}{)}\PYG{p}{:}
\PYG{o}{\PYGZgt{}\PYGZgt{}}     \PYG{k}{if} \PYG{n}{k} \PYG{o}{==} \PYG{l+m+mi}{0}\PYG{p}{:}
\PYG{o}{\PYGZgt{}\PYGZgt{}}         \PYG{k}{return} \PYG{l+m+mf}{0.0}
\PYG{o}{\PYGZgt{}\PYGZgt{}}     \PYG{k}{else}\PYG{p}{:}
\PYG{o}{\PYGZgt{}\PYGZgt{}}         \PYG{k}{return}  \PYG{n}{A}\PYG{o}{*}\PYG{p}{(}\PYG{n}{k}\PYG{o}{/}\PYG{n}{k0}\PYG{p}{)}\PYG{o}{*}\PYG{o}{*}\PYG{p}{(}\PYG{n}{s}\PYG{p}{)}
\PYG{o}{\PYGZgt{}\PYGZgt{}}
\PYG{o}{\PYGZgt{}\PYGZgt{}} \PYG{n}{sv}\PYG{o}{.}\PYG{n}{sim\PYGZus{}vis}\PYG{p}{(}\PYG{n}{infits}\PYG{p}{,} \PYG{n}{outfits}\PYG{p}{,} \PYG{n}{skysimtype} \PYG{o}{=} \PYG{l+s+s1}{\PYGZsq{}}\PYG{l+s+s1}{3D}\PYG{l+s+s1}{\PYGZsq{}}\PYG{p}{,}
                    \PYG{n}{seed} \PYG{o}{=} \PYG{l+m+mi}{5}\PYG{p}{,} \PYG{n}{psfunc} \PYG{o}{=} \PYG{n}{mypkfunc}\PYG{p}{)}
\end{sphinxVerbatim}
\end{quote}
\subsubsection*{Notes}
\begin{itemize}
\item {} 
\sphinxAtStartPar
\sphinxcode{\sphinxupquote{APS}} function generates \sphinxcode{\sphinxupquote{2D sky signal}} at centered frequency \(\nu_c\).

\item {} 
\sphinxAtStartPar
\sphinxcode{\sphinxupquote{PS}} function generates \sphinxcode{\sphinxupquote{3D sky signal}} for all frequency channels available in the \sphinxstylestrong{fits file}.

\end{itemize}


\bigskip\hrule\bigskip



\subsection{Important point}
\label{\detokenize{simvis:important-point}}
\sphinxAtStartPar
Before you can simulate, you need to first set some parameters related to simulation.
\begin{itemize}
\item {} 
\sphinxAtStartPar
\sphinxcode{\sphinxupquote{nside}} : Healpix parameter that sets the resolution.

\item {} 
\sphinxAtStartPar
\sphinxcode{\sphinxupquote{Nrea}} : No of realizations of sky (\(N_{\rm realizations}\)).

\item {} 
\sphinxAtStartPar
\sphinxcode{\sphinxupquote{chunk\_size}} : Number of baselines for which visibilities are being simulated at one step.

\end{itemize}

\sphinxAtStartPar
\sphinxcode{\sphinxupquote{chunk\_size}} is a parameter which can be varied as per system requirements. If you have higher ram say 32 GB use it as 300. For 16 you can use 200 and and for less use 100. This number is for \sphinxcode{\sphinxupquote{nside}} = 512 and \sphinxcode{\sphinxupquote{Nrea}} = 100. As in a part of calculations a matrix comes which has the dimensions of  \(6\text{nside}^2 \times \text{chunk\_size}\) with each entry represented by 128 bit complex number. So it takes memory to store it in RAM. This should not be more than RAM capacity. If you have more RAM you can’t increase the \sphinxcode{\sphinxupquote{chunk\_size}} indefinitely as matrix operations will take much time to do the computation. I found, for \sphinxcode{\sphinxupquote{nside}} = 512 and \sphinxcode{\sphinxupquote{Nrea}} = 100 , \sphinxcode{\sphinxupquote{chunk\_size = 300\sphinxhyphen{}400}} (for 32 GB device) is efficient.

\sphinxAtStartPar
\sphinxstylestrong{To set them, follow this :}

\begin{sphinxVerbatim}[commandchars=\\\{\}]
\PYG{o}{\PYGZgt{}\PYGZgt{}} \PYG{k+kn}{import} \PYG{n+nn}{simvis} \PYG{k}{as} \PYG{n+nn}{sv}
\PYG{o}{\PYGZgt{}\PYGZgt{}} \PYG{k+kn}{import} \PYG{n+nn}{builtins} \PYG{k}{as} \PYG{n+nn}{blt}
\PYG{o}{\PYGZgt{}\PYGZgt{}}
\PYG{o}{\PYGZgt{}\PYGZgt{}} \PYG{c+c1}{\PYGZsh{} Set simulation parameters}
\PYG{o}{\PYGZgt{}\PYGZgt{}} \PYG{n}{blt}\PYG{o}{.}\PYG{n}{nside} \PYG{o}{=} \PYG{l+m+mi}{512}        \PYG{c+c1}{\PYGZsh{} nside parameter}
\PYG{o}{\PYGZgt{}\PYGZgt{}} \PYG{n}{blt}\PYG{o}{.}\PYG{n}{Nrea}  \PYG{o}{=} \PYG{l+m+mi}{100}        \PYG{c+c1}{\PYGZsh{} no of realizations}
\PYG{o}{\PYGZgt{}\PYGZgt{}} \PYG{n}{blt}\PYG{o}{.}\PYG{n}{chunk\PYGZus{}size} \PYG{o}{=} \PYG{l+m+mi}{400}   \PYG{c+c1}{\PYGZsh{} chunk\PYGZus{}size, no of baselines in a single step}
\end{sphinxVerbatim}


\section{Overview}
\label{\detokenize{simvis:overview}}\begin{itemize}
\item {} 
\sphinxAtStartPar
First we generate the \sphinxstylestrong{Gaussian Random Field} \(T(\hat{\mathbf{n}},\nu)\)  (temperature field in unit of mK) from  \sphinxstylestrong{Angular Power Spectrum} or \sphinxstylestrong{Power Spectrum}.

\item {} 
\sphinxAtStartPar
Calculate the \sphinxstylestrong{primary beam pattern} \({A}(\Delta\hat{\mathbf{n}},\nu)\) for the telescope.

\item {} 
\sphinxAtStartPar
Calculate the components of \(\hat{\mathbf{n}}\) along the basis vectors \(\hat{\mathbf{e}}_1,\hat{\mathbf{e}}_2,\hat{\mathbf{e}}_3\). This will help us to determine the dot product \(\vec{\textbf{U}}\cdot\Delta\hat{\mathbf{n}}\).

\item {} 
\sphinxAtStartPar
Now we calculate the \sphinxstylestrong{phase factor} \(\textbf{exp}\left[{2\pi \text{i}\vec{\textbf{U}}\cdot\Delta\hat{\mathbf{n}}}\right]\).

\item {} 
\sphinxAtStartPar
Multiply the \sphinxstylestrong{GRF and PB} first and then multiply by phase factor, sum over the pixels. The whole thing is \sphinxtitleref{done by matrix multiplication.}

\item {} 
\sphinxAtStartPar
The visibility is simulated.

\end{itemize}

\clearpage


\section{Module Functions}
\label{\detokenize{simvis:module-functions}}\index{aps() (in module simvis)@\spxentry{aps()}\spxextra{in module simvis}}

\begin{fulllineitems}
\phantomsection\label{\detokenize{simvis:simvis.aps}}
\pysigstartsignatures
\pysiglinewithargsret{\sphinxcode{\sphinxupquote{simvis.}}\sphinxbfcode{\sphinxupquote{aps}}}{\sphinxparam{\DUrole{n}{l}}}{}
\pysigstopsignatures
\sphinxAtStartPar
This is the input model Angular Power Spectrum(APS) function.
\begin{equation*}
\begin{split}C_{\ell}=A\, \left(\frac{\ell}{\ell_0}\right)^{\beta}\end{split}
\end{equation*}
\sphinxAtStartPar
A defines the amplitude of the APS is in the unit of \(\text{mK}^2\). \(\beta\) here is the power\sphinxhyphen{}law index. \(C_{\ell}\) has the unit of  (\(\text{mK}^2\)). By default set to \(A = 100\;\text{mK}^2\), \(\ell_0=1\), \(\beta\) = \sphinxhyphen{}2.
\begin{quote}\begin{description}
\sphinxlineitem{Parameters}
\sphinxAtStartPar
\sphinxstyleliteralstrong{\sphinxupquote{l}} (\sphinxstyleliteralemphasis{\sphinxupquote{float}}\sphinxstyleliteralemphasis{\sphinxupquote{ or }}\sphinxstyleliteralemphasis{\sphinxupquote{np.ndarray}}) \textendash{} The angular multipole \({\ell}\) value(s).

\sphinxlineitem{Returns}
\sphinxAtStartPar
The angular power spectrum value for the desired multipole(s) \({\ell}\).

\sphinxlineitem{Return type}
\sphinxAtStartPar
float or np.ndarray

\end{description}\end{quote}

\end{fulllineitems}

\index{uaps() (in module simvis)@\spxentry{uaps()}\spxextra{in module simvis}}

\begin{fulllineitems}
\phantomsection\label{\detokenize{simvis:simvis.uaps}}
\pysigstartsignatures
\pysiglinewithargsret{\sphinxcode{\sphinxupquote{simvis.}}\sphinxbfcode{\sphinxupquote{uaps}}}{\sphinxparam{\DUrole{n}{l}}}{}
\pysigstopsignatures
\sphinxAtStartPar
This is the Unit Angular Power Spectrum function.
\begin{quote}\begin{description}
\sphinxlineitem{Parameters}
\sphinxAtStartPar
\sphinxstyleliteralstrong{\sphinxupquote{l}} (\sphinxstyleliteralemphasis{\sphinxupquote{float}}\sphinxstyleliteralemphasis{\sphinxupquote{ or }}\sphinxstyleliteralemphasis{\sphinxupquote{np.ndarray}}) \textendash{} The angular multipole \({\ell}\) value(s).

\sphinxlineitem{Returns}
\sphinxAtStartPar
1.0

\sphinxlineitem{Return type}
\sphinxAtStartPar
float

\end{description}\end{quote}

\end{fulllineitems}

\index{grf() (in module simvis)@\spxentry{grf()}\spxextra{in module simvis}}

\begin{fulllineitems}
\phantomsection\label{\detokenize{simvis:simvis.grf}}
\pysigstartsignatures
\pysiglinewithargsret{\sphinxcode{\sphinxupquote{simvis.}}\sphinxbfcode{\sphinxupquote{grf}}}{\sphinxparam{\DUrole{n}{nside}}\sphinxparamcomma \sphinxparam{\DUrole{n}{lmax}}\sphinxparamcomma \sphinxparam{\DUrole{n}{apsfunc}}\sphinxparamcomma \sphinxparam{\DUrole{n}{seed}}}{}
\pysigstopsignatures
\sphinxAtStartPar
Generates  GRF values  for every pixel using the input APS(apsfunc).The GRF is monopole subtracted.
\begin{quote}\begin{description}
\sphinxlineitem{Parameters}\begin{itemize}
\item {} 
\sphinxAtStartPar
\sphinxstyleliteralstrong{\sphinxupquote{nside}} (\sphinxstyleliteralemphasis{\sphinxupquote{int}}) \textendash{} Healpix parameter, usually power of 2.

\item {} 
\sphinxAtStartPar
\sphinxstyleliteralstrong{\sphinxupquote{lmax}} (\sphinxstyleliteralemphasis{\sphinxupquote{int}}) \textendash{} Maximum value of angular multipole.(\(\ell\))

\item {} 
\sphinxAtStartPar
\sphinxstyleliteralstrong{\sphinxupquote{apsfunc}} (\sphinxstyleliteralemphasis{\sphinxupquote{function}}) \textendash{} The input APS function.

\item {} 
\sphinxAtStartPar
\sphinxstyleliteralstrong{\sphinxupquote{seed}} (\sphinxstyleliteralemphasis{\sphinxupquote{int}}\sphinxstyleliteralemphasis{\sphinxupquote{ or }}\sphinxstyleliteralemphasis{\sphinxupquote{float}}) \textendash{} A seed value, which can be provied to get different realizations for same APS.
If you pass the seed value, then the output GRF is a fixed one.

\end{itemize}

\sphinxlineitem{Returns}
\sphinxAtStartPar
Contains the GRF value for each pixels. Shape (\(12\ {nside}^2,\)). Returned values are in  \(\text{mK}\) unit.

\sphinxlineitem{Return type}
\sphinxAtStartPar
np.ndarray

\end{description}\end{quote}

\clearpage

\end{fulllineitems}

\index{ps() (in module simvis)@\spxentry{ps()}\spxextra{in module simvis}}

\begin{fulllineitems}
\phantomsection\label{\detokenize{simvis:simvis.ps}}
\pysigstartsignatures
\pysiglinewithargsret{\sphinxcode{\sphinxupquote{simvis.}}\sphinxbfcode{\sphinxupquote{ps}}}{\sphinxparam{\DUrole{n}{k}}}{}
\pysigstopsignatures
\sphinxAtStartPar
This is the input model Power Spectrum function, which has the form \(A \left(\dfrac{k}{k_0}\right)^{s}\), with \(A=1, k_0=1, s=-2\) (Default set).
\(A\) has the unit of \(\text{mK}^2\;\text{Mpc}^3\). \(k, k_0\) are in \(\text{Mpc}^{-1}\) unit.
\begin{quote}\begin{description}
\sphinxlineitem{Parameters}
\sphinxAtStartPar
\sphinxstyleliteralstrong{\sphinxupquote{k}} (\sphinxstyleliteralemphasis{\sphinxupquote{float}}\sphinxstyleliteralemphasis{\sphinxupquote{ or }}\sphinxstyleliteralemphasis{\sphinxupquote{np.ndarray}}) \textendash{} \(k = |\vec{\mathbf{k}}|\), the comoving wave number.

\sphinxlineitem{Returns}
\sphinxAtStartPar
The value of Power Spectrum \(P(k)\)  at comoving wave number \(k\).

\sphinxlineitem{Return type}
\sphinxAtStartPar
float or np.ndarray

\end{description}\end{quote}
\subsubsection*{Notes}

\sphinxAtStartPar
\sphinxstylestrong{Returns 0.0 if} \(k=0\).

\end{fulllineitems}

\index{sky3dgrf() (in module simvis)@\spxentry{sky3dgrf()}\spxextra{in module simvis}}

\begin{fulllineitems}
\phantomsection\label{\detokenize{simvis:simvis.sky3dgrf}}
\pysigstartsignatures
\pysiglinewithargsret{\sphinxcode{\sphinxupquote{simvis.}}\sphinxbfcode{\sphinxupquote{sky3dgrf}}}{\sphinxparam{\DUrole{n}{nside}}\sphinxparamcomma \sphinxparam{\DUrole{n}{seed}}\sphinxparamcomma \sphinxparam{\DUrole{n}{psfunc=\textless{}function ps\textgreater{}}}}{}
\pysigstopsignatures
\sphinxAtStartPar
This function simulates the 3D sky signal for a given \sphinxcode{\sphinxupquote{nside}}, \sphinxcode{\sphinxupquote{Power Spectrum function}}, \sphinxcode{\sphinxupquote{seed}} value across all the frequency channels taken from  input \sphinxcode{\sphinxupquote{UVFITS}} file.
\begin{quote}\begin{description}
\sphinxlineitem{Parameters}\begin{itemize}
\item {} 
\sphinxAtStartPar
\sphinxstyleliteralstrong{\sphinxupquote{nside}} (\sphinxstyleliteralemphasis{\sphinxupquote{int}}) \textendash{} Healpix parameter (usually power of 2)

\item {} 
\sphinxAtStartPar
\sphinxstyleliteralstrong{\sphinxupquote{seed}} (\sphinxstyleliteralemphasis{\sphinxupquote{Positive int}}\sphinxstyleliteralemphasis{\sphinxupquote{ or }}\sphinxstyleliteralemphasis{\sphinxupquote{None}}) \textendash{} seed value for generating the \(a_{\ell m}\)\textquotesingle{}s. Different seed value will result in different \(a_{\ell m}\) \textquotesingle{}s, but for a particular seed value that will be fixed.

\item {} 
\sphinxAtStartPar
\sphinxstyleliteralstrong{\sphinxupquote{psfunc}} (\sphinxstyleliteralemphasis{\sphinxupquote{function}}) \textendash{} The input Power Spectrum function. By default it is \(\texttt{ps}\) which has the form \(A \left(\dfrac{k}{k_0}\right)^{s}\), with \(A=1, k_0=1, s=-2\).
You can define your own power spectrum function make sure that it has only  \(k\) dependence.

\end{itemize}

\sphinxlineitem{Returns}
\sphinxAtStartPar
2D array has the shape \(N_c \times \text{npix}\). Each row data corresponds to GRF at a particular frequency. npix is the number of pixels. To get the value of npix see notes section below.(usually it is \(12\text{nside}^2\)).

\sphinxlineitem{Return type}
\sphinxAtStartPar
Array type

\end{description}\end{quote}

\clearpage
\subsubsection*{Notes}

\sphinxAtStartPar
By default \(\ell_{max} = 3*\text{nside} -1\).
To extract gaussian random field value \(T(\hat{\mathbf{n}}, \nu)\) for a given channel no or no\textquotesingle{}s (indexing starts from 0).

\begin{sphinxVerbatim}[commandchars=\\\{\}]
\PYG{o}{\PYGZgt{}\PYGZgt{}} \PYG{k+kn}{import} \PYG{n+nn}{healpy} \PYG{k}{as} \PYG{n+nn}{hp}
\PYG{o}{\PYGZgt{}\PYGZgt{}} \PYG{n}{nside}\PYG{p}{,}\PYG{n}{seed} \PYG{o}{=} \PYG{l+m+mi}{16}\PYG{p}{,} \PYG{k+kc}{None}
\PYG{o}{\PYGZgt{}\PYGZgt{}} \PYG{n}{npix} \PYG{o}{=} \PYG{n}{hp}\PYG{o}{.}\PYG{n}{nside2npix}\PYG{p}{(}\PYG{n}{nside}\PYG{p}{)}
\PYG{o}{\PYGZgt{}\PYGZgt{}} \PYG{n+nb}{print}\PYG{p}{(}\PYG{n}{npix}\PYG{p}{)}
\PYG{o}{\PYGZgt{}\PYGZgt{}} \PYG{n}{channel} \PYG{o}{=} \PYG{l+m+mi}{10}
\PYG{o}{\PYGZgt{}\PYGZgt{}} \PYG{c+c1}{\PYGZsh{} for channels from 1 to 10}
\PYG{o}{\PYGZgt{}\PYGZgt{}} \PYG{n}{channel} \PYG{o}{=} \PYG{p}{[} \PYG{n}{i} \PYG{k}{for} \PYG{n}{i} \PYG{o+ow}{in} \PYG{n+nb}{range}\PYG{p}{(}\PYG{l+m+mi}{10}\PYG{p}{)}\PYG{p}{]}
\PYG{o}{\PYGZgt{}\PYGZgt{}} \PYG{n}{grf} \PYG{o}{=} \PYG{n}{sky3dgrf}\PYG{p}{(}\PYG{n}{nside}\PYG{p}{,}\PYG{n}{seed}\PYG{p}{)} \PYG{c+c1}{\PYGZsh{} generate the grf}
\PYG{o}{\PYGZgt{}\PYGZgt{}} \PYG{n}{grf}\PYG{p}{[}\PYG{n}{channel}\PYG{p}{]}             \PYG{c+c1}{\PYGZsh{} extract the values}
\end{sphinxVerbatim}

\sphinxAtStartPar
GRF across all channels has no monopole, as \(a_{0,0}\) \sphinxstylestrong{is set zero.}


\bigskip\hrule\bigskip


\sphinxAtStartPar
If you want to give your own power spectrum function (make sure it is vectorized)

\begin{sphinxVerbatim}[commandchars=\\\{\}]
\PYG{o}{\PYGZgt{}\PYGZgt{}} \PYG{k+kn}{import} \PYG{n+nn}{numpy} \PYG{k}{as} \PYG{n+nn}{np}
\PYG{o}{\PYGZgt{}\PYGZgt{}} \PYG{k+kn}{import} \PYG{n+nn}{builtins} \PYG{k}{as} \PYG{n+nn}{blt}
\PYG{o}{\PYGZgt{}\PYGZgt{}} \PYG{k+kn}{import} \PYG{n+nn}{simvis} \PYG{k}{as} \PYG{n+nn}{sv}
\PYG{o}{\PYGZgt{}\PYGZgt{}} \PYG{c+c1}{\PYGZsh{} make sure all are float.}
\PYG{o}{\PYGZgt{}\PYGZgt{}} \PYG{n}{blt}\PYG{o}{.}\PYG{n}{A}  \PYG{o}{=}  \PYG{l+m+mf}{1.0}
\PYG{o}{\PYGZgt{}\PYGZgt{}} \PYG{n}{blt}\PYG{o}{.}\PYG{n}{s}  \PYG{o}{=} \PYG{o}{\PYGZhy{}}\PYG{l+m+mf}{3.0}
\PYG{o}{\PYGZgt{}\PYGZgt{}} \PYG{n}{blt}\PYG{o}{.}\PYG{n}{k0} \PYG{o}{=}  \PYG{l+m+mf}{1.0}
\PYG{o}{\PYGZgt{}\PYGZgt{}} \PYG{n+nd}{@np}\PYG{o}{.}\PYG{n}{vectorize}
\PYG{o}{\PYGZgt{}\PYGZgt{}} \PYG{k}{def} \PYG{n+nf}{mypkfunc}\PYG{p}{(}\PYG{n}{k}\PYG{p}{)}\PYG{p}{:}
\PYG{o}{\PYGZgt{}\PYGZgt{}}     \PYG{k}{if} \PYG{n}{k} \PYG{o}{==} \PYG{l+m+mi}{0}\PYG{p}{:}
\PYG{o}{\PYGZgt{}\PYGZgt{}}         \PYG{k}{return} \PYG{l+m+mf}{0.0}
\PYG{o}{\PYGZgt{}\PYGZgt{}}     \PYG{k}{else}\PYG{p}{:}
\PYG{o}{\PYGZgt{}\PYGZgt{}}         \PYG{k}{return}  \PYG{n}{A}\PYG{o}{*}\PYG{p}{(}\PYG{n}{k}\PYG{o}{/}\PYG{n}{k0}\PYG{p}{)}\PYG{o}{*}\PYG{o}{*}\PYG{p}{(}\PYG{n}{s}\PYG{p}{)}
\PYG{o}{\PYGZgt{}\PYGZgt{}}
\PYG{o}{\PYGZgt{}\PYGZgt{}} \PYG{n}{grf} \PYG{o}{=} \PYG{n}{sv}\PYG{o}{.}\PYG{n}{sky3dgrf}\PYG{p}{(}\PYG{n}{nside}\PYG{p}{,} \PYG{n}{seed}\PYG{p}{,} \PYG{n}{psfunc} \PYG{o}{=} \PYG{n}{mypkfunc}\PYG{p}{)}  \PYG{c+c1}{\PYGZsh{} generate the grf}
\PYG{o}{\PYGZgt{}\PYGZgt{}}
\end{sphinxVerbatim}

\clearpage

\end{fulllineitems}

\index{allskysim() (in module simvis)@\spxentry{allskysim()}\spxextra{in module simvis}}

\begin{fulllineitems}
\phantomsection\label{\detokenize{simvis:simvis.allskysim}}
\pysigstartsignatures
\pysiglinewithargsret{\sphinxcode{\sphinxupquote{simvis.}}\sphinxbfcode{\sphinxupquote{allskysim}}}{\sphinxparam{\DUrole{n}{nside}}\sphinxparamcomma \sphinxparam{\DUrole{n}{seed}}\sphinxparamcomma \sphinxparam{\DUrole{n}{skysimtype}}\sphinxparamcomma \sphinxparam{\DUrole{n}{apsfunc=\textless{}function aps\textgreater{}}}\sphinxparamcomma \sphinxparam{\DUrole{n}{psfunc=\textless{}function ps\textgreater{}}}}{}
\pysigstopsignatures
\sphinxAtStartPar
This function creates all sky maps accordingly.
\begin{quote}\begin{description}
\sphinxlineitem{Parameters}\begin{itemize}
\item {} 
\sphinxAtStartPar
\sphinxstyleliteralstrong{\sphinxupquote{nside}} (\sphinxstyleliteralemphasis{\sphinxupquote{int}}) \textendash{} Healpix parameter.Usually power of 2.

\item {} 
\sphinxAtStartPar
\sphinxstyleliteralstrong{\sphinxupquote{seed}} (\sphinxstyleliteralemphasis{\sphinxupquote{Positive int}}\sphinxstyleliteralemphasis{\sphinxupquote{ or }}\sphinxstyleliteralemphasis{\sphinxupquote{None}}) \textendash{} If it’s 2D map then in each realization the seed gets increased by 1.

\item {} 
\sphinxAtStartPar
\sphinxstyleliteralstrong{\sphinxupquote{skysimtype}} (\sphinxstyleliteralemphasis{\sphinxupquote{str}}) \textendash{} Either ‘2D’ or ‘3D’. If given ‘2D’ then creates multiple realizations. If given ‘3D’ then makes a 3D sky map.

\item {} 
\sphinxAtStartPar
\sphinxstyleliteralstrong{\sphinxupquote{apsfunc}} (\sphinxstyleliteralemphasis{\sphinxupquote{function}}) \textendash{} The input Angular Power Spectrum function. By default it is \(\texttt{aps}\) which has the form \(C_{\ell}=A\, \left(\frac{\ell}{\ell_0}\right)^{\beta}\) with \(A = 100\;\text{mK}^2\), \(\ell_0=1\), \(\beta\) = \sphinxhyphen{}2.

\item {} 
\sphinxAtStartPar
\sphinxstyleliteralstrong{\sphinxupquote{psfunc}} (\sphinxstyleliteralemphasis{\sphinxupquote{function}}) \textendash{} The input Power Spectrum function. By default it is \(\texttt{ps}\) which has the form \(A \left(\dfrac{k}{k_0}\right)^{s}\), with \(A=1, k_0=1, s=-2\).

\end{itemize}

\sphinxlineitem{Returns}
\sphinxAtStartPar
Surface brightness temperature field map (multiple realizations or across all frequencies) in units of \(\text{mK}\).

\sphinxlineitem{Return type}
\sphinxAtStartPar
np.ndarray

\end{description}\end{quote}

\end{fulllineitems}

\index{beammwa() (in module simvis)@\spxentry{beammwa()}\spxextra{in module simvis}}

\begin{fulllineitems}
\phantomsection\label{\detokenize{simvis:simvis.beammwa}}
\pysigstartsignatures
\pysiglinewithargsret{\sphinxcode{\sphinxupquote{simvis.}}\sphinxbfcode{\sphinxupquote{beammwa}}}{\sphinxparam{\DUrole{n}{nu}}\sphinxparamcomma \sphinxparam{\DUrole{n}{ne1}}\sphinxparamcomma \sphinxparam{\DUrole{n}{ne2}}}{}
\pysigstopsignatures
\sphinxAtStartPar
We model primary beam function of \sphinxstylestrong{MWA telescope} as a product of two sinc square function given as :
\begin{equation*}
\begin{split}\mathcal{A}(\hat{\mathbf{n}},\nu) = \text{sinc}^2\left[\frac{{\pi}b{\nu}\ \Delta\hat{\mathbf{n}}\cdot\hat{\mathbf{e}}_1 (\alpha_{p})}{c}\right]\;\text{sinc}^2\left[\frac{{\pi}b{\nu}\ \Delta\hat{\mathbf{n}}\cdot\hat{\mathbf{e}}_2 (\alpha_{p})}{c}\right]\end{split}
\end{equation*}
\sphinxAtStartPar
Where :
\begin{equation*}
\begin{split}\Delta\hat{\mathbf{n}}= \hat{\mathbf{n}}- \hat{\mathbf{p}}\end{split}
\end{equation*}
\sphinxAtStartPar
\(\hat{\mathbf{p}}\) is the pointing direction of the telescope. We assume that telescope is zenith pointed (\(\hat{\mathbf{p}}\) = \(\hat{\mathbf{e}}_3\)).
\(\hat{\mathbf{n}}\) refers to different directions on the sky (i.e. celestial sphere).
\begin{quote}\begin{description}
\sphinxlineitem{Parameters}\begin{itemize}
\item {} 
\sphinxAtStartPar
\sphinxstyleliteralstrong{\sphinxupquote{nu}} (\sphinxstyleliteralemphasis{\sphinxupquote{float}}) \textendash{} Observing frequency \(\nu\) in MHz unit.

\item {} 
\sphinxAtStartPar
\sphinxstyleliteralstrong{\sphinxupquote{ne1}} (\sphinxstyleliteralemphasis{\sphinxupquote{float}}\sphinxstyleliteralemphasis{\sphinxupquote{ or }}\sphinxstyleliteralemphasis{\sphinxupquote{array}}) \textendash{} The dot product given by \(\hat{\mathbf{n}}\cdot\hat{\mathbf{e}}_ 1\).

\item {} 
\sphinxAtStartPar
\sphinxstyleliteralstrong{\sphinxupquote{ne2}} (\sphinxstyleliteralemphasis{\sphinxupquote{float}}\sphinxstyleliteralemphasis{\sphinxupquote{ or }}\sphinxstyleliteralemphasis{\sphinxupquote{array}}) \textendash{} The dot product given by \(\hat{\mathbf{n}}\cdot\hat{\mathbf{e}}_ 2\).

\end{itemize}

\sphinxlineitem{Returns}
\sphinxAtStartPar
The primary beam pattern value for square aperture telescope (MWA).

\sphinxlineitem{Return type}
\sphinxAtStartPar
float or np.ndarray

\end{description}\end{quote}

\clearpage

\end{fulllineitems}

\index{pbgen() (in module simvis)@\spxentry{pbgen()}\spxextra{in module simvis}}

\begin{fulllineitems}
\phantomsection\label{\detokenize{simvis:simvis.pbgen}}
\pysigstartsignatures
\pysiglinewithargsret{\sphinxcode{\sphinxupquote{simvis.}}\sphinxbfcode{\sphinxupquote{pbgen}}}{\sphinxparam{\DUrole{n}{nside}}\sphinxparamcomma \sphinxparam{\DUrole{n}{ra\_ptg}}\sphinxparamcomma \sphinxparam{\DUrole{n}{dec\_ptg}}\sphinxparamcomma \sphinxparam{\DUrole{n}{nu}}\sphinxparamcomma \sphinxparam{\DUrole{n}{pbfunction}}}{}
\pysigstopsignatures
\sphinxAtStartPar
This generates the primary beam for MWA telescope (zenith pointed), given the RA,DEC and frequency.
\begin{quote}\begin{description}
\sphinxlineitem{Parameters}\begin{itemize}
\item {} 
\sphinxAtStartPar
\sphinxstyleliteralstrong{\sphinxupquote{nside}} (\sphinxstyleliteralemphasis{\sphinxupquote{int}}) \textendash{} Healpix parameter.Usually power of 2.

\item {} 
\sphinxAtStartPar
\sphinxstyleliteralstrong{\sphinxupquote{ra\_ptg}} (\sphinxstyleliteralemphasis{\sphinxupquote{int}}\sphinxstyleliteralemphasis{\sphinxupquote{ or }}\sphinxstyleliteralemphasis{\sphinxupquote{float}}) \textendash{} RA of the pointing direction.(In degrees)(zenith pointed).

\item {} 
\sphinxAtStartPar
\sphinxstyleliteralstrong{\sphinxupquote{dec\_ptg}} (\sphinxstyleliteralemphasis{\sphinxupquote{int}}\sphinxstyleliteralemphasis{\sphinxupquote{ or }}\sphinxstyleliteralemphasis{\sphinxupquote{float}}) \textendash{} DEC of the pointing direction.(In degrees)(zenith pointed).

\item {} 
\sphinxAtStartPar
\sphinxstyleliteralstrong{\sphinxupquote{nu}} (\sphinxstyleliteralemphasis{\sphinxupquote{float}}) \textendash{} Observing frequency \(\nu\) in \(\mathrm{MHz}\) unit.

\item {} 
\sphinxAtStartPar
\sphinxstyleliteralstrong{\sphinxupquote{pbfunction}} (\sphinxstyleliteralemphasis{\sphinxupquote{function}}) \textendash{} The primary beam (PB) pattern of MWA which is product of two \(\text{sinc}^2\) functions.

\end{itemize}

\sphinxlineitem{Returns}
\sphinxAtStartPar
\sphinxstylestrong{PB} \textendash{} The value of primary beam of the MWA telescope (beam pattern value at each pixel which quantifies the reponse of telescope along different directions). Shape \((12\ {nside}^2, )\).

\sphinxlineitem{Return type}
\sphinxAtStartPar
np.ndarray

\end{description}\end{quote}

\end{fulllineitems}

\index{hat\_n() (in module simvis)@\spxentry{hat\_n()}\spxextra{in module simvis}}

\begin{fulllineitems}
\phantomsection\label{\detokenize{simvis:simvis.hat_n}}
\pysigstartsignatures
\pysiglinewithargsret{\sphinxcode{\sphinxupquote{simvis.}}\sphinxbfcode{\sphinxupquote{hat\_n}}}{\sphinxparam{\DUrole{n}{nside}}\sphinxparamcomma \sphinxparam{\DUrole{n}{ipix}}}{}
\pysigstopsignatures
\sphinxAtStartPar
This calculates the \(xyz\) components of \(\hat{\bf n}\) for all the given pixels contained in the array ipix.
\begin{quote}\begin{description}
\sphinxlineitem{Parameters}\begin{itemize}
\item {} 
\sphinxAtStartPar
\sphinxstyleliteralstrong{\sphinxupquote{nside}} (\sphinxstyleliteralemphasis{\sphinxupquote{int}}) \textendash{} Healpix parameter.Usually power of 2.

\item {} 
\sphinxAtStartPar
\sphinxstyleliteralstrong{\sphinxupquote{ipix}} (\sphinxstyleliteralemphasis{\sphinxupquote{np.ndarray}}) \textendash{} Pixel array for which components of \(\hat{\bf n}\) have to be calculated.

\end{itemize}

\sphinxlineitem{Returns}
\sphinxAtStartPar
\sphinxstylestrong{hatn} \textendash{} Containg the \(xyz\) comps of \(\hat{\bf n}\). \(N_{pix} \times 3\).

\sphinxlineitem{Return type}
\sphinxAtStartPar
np.ndarray

\end{description}\end{quote}

\end{fulllineitems}

\index{dot\_cal\_superfast() (in module simvis)@\spxentry{dot\_cal\_superfast()}\spxextra{in module simvis}}

\begin{fulllineitems}
\phantomsection\label{\detokenize{simvis:simvis.dot_cal_superfast}}
\pysigstartsignatures
\pysiglinewithargsret{\sphinxcode{\sphinxupquote{simvis.}}\sphinxbfcode{\sphinxupquote{dot\_cal\_superfast}}}{\sphinxparam{\DUrole{n}{nside}}\sphinxparamcomma \sphinxparam{\DUrole{n}{ra\_ptg}}\sphinxparamcomma \sphinxparam{\DUrole{n}{dec\_ptg}}\sphinxparamcomma \sphinxparam{\DUrole{n}{ipix}}}{}
\pysigstopsignatures
\sphinxAtStartPar
This calculates components of \(\Delta\hat{\bf n}\) for all the given pixels contained in the array ipix along the basis vectors \(\hat{\mathbf{e}}_1,\hat{\mathbf{e}}_2,\hat{\mathbf{e}}_3\) given by the RA,DEC.
\begin{quote}\begin{description}
\sphinxlineitem{Parameters}\begin{itemize}
\item {} 
\sphinxAtStartPar
\sphinxstyleliteralstrong{\sphinxupquote{nside}} (\sphinxstyleliteralemphasis{\sphinxupquote{int}}) \textendash{} Healpix parameter. Usually power of 2.

\item {} 
\sphinxAtStartPar
\sphinxstyleliteralstrong{\sphinxupquote{ra\_ptg}} (\sphinxstyleliteralemphasis{\sphinxupquote{int}}\sphinxstyleliteralemphasis{\sphinxupquote{ or }}\sphinxstyleliteralemphasis{\sphinxupquote{float}}) \textendash{} RA of the pointing direction.(In degrees)

\item {} 
\sphinxAtStartPar
\sphinxstyleliteralstrong{\sphinxupquote{dec\_ptg}} (\sphinxstyleliteralemphasis{\sphinxupquote{int}}\sphinxstyleliteralemphasis{\sphinxupquote{ or }}\sphinxstyleliteralemphasis{\sphinxupquote{float}}) \textendash{} DEC of the pointing direction.(In degrees)

\item {} 
\sphinxAtStartPar
\sphinxstyleliteralstrong{\sphinxupquote{ipix}} (\sphinxstyleliteralemphasis{\sphinxupquote{np.ndarray}}) \textendash{} Pixel array for which dot products have to be calculated.

\end{itemize}

\sphinxlineitem{Returns}
\sphinxAtStartPar
\sphinxstylestrong{dot\_products} \textendash{} Contains the \(xyz\) comps for those pixels. Shape \(N_{pix} \times 3\).

\sphinxlineitem{Return type}
\sphinxAtStartPar
np.ndarray

\end{description}\end{quote}

\end{fulllineitems}

\index{calculate\_phase() (in module simvis)@\spxentry{calculate\_phase()}\spxextra{in module simvis}}

\begin{fulllineitems}
\phantomsection\label{\detokenize{simvis:simvis.calculate_phase}}
\pysigstartsignatures
\pysiglinewithargsret{\sphinxcode{\sphinxupquote{simvis.}}\sphinxbfcode{\sphinxupquote{calculate\_phase}}}{\sphinxparam{\DUrole{n}{dot\_product}}\sphinxparamcomma \sphinxparam{\DUrole{n}{bl}}}{}
\pysigstopsignatures
\sphinxAtStartPar
Given the baselines it calculates the phase factor for all the pixels given in dot\_product array.
\begin{quote}\begin{description}
\sphinxlineitem{Parameters}\begin{itemize}
\item {} 
\sphinxAtStartPar
\sphinxstyleliteralstrong{\sphinxupquote{dot\_product}} (\sphinxstyleliteralemphasis{\sphinxupquote{np.ndarray}}) \textendash{} Contains the  components for all the given those pixels along the basis vectors \(\hat{\mathbf{e}}_1,\hat{\mathbf{e}}_2,\hat{\mathbf{e}}_3\). Shape is \(N_{pix} \times 3\).

\item {} 
\sphinxAtStartPar
\sphinxstyleliteralstrong{\sphinxupquote{bl}} (\sphinxstyleliteralemphasis{\sphinxupquote{np.ndarray}}) \textendash{} The array conatins the components of the baselines. (defined by basis vectors \(\hat{\mathbf{e}}_1,\hat{\mathbf{e}}_2,\hat{\mathbf{e}}_3\)).
Shape \(N_{Baselines}\times 3\).

\end{itemize}

\sphinxlineitem{Returns}
\sphinxAtStartPar
The phase factor. Shape \(N_{pix} \times  N_{Baselines}\).

\sphinxlineitem{Return type}
\sphinxAtStartPar
np.ndarray

\end{description}\end{quote}

\end{fulllineitems}

\index{visgen\_mwa\_multi() (in module simvis)@\spxentry{visgen\_mwa\_multi()}\spxextra{in module simvis}}

\begin{fulllineitems}
\phantomsection\label{\detokenize{simvis:simvis.visgen_mwa_multi}}
\pysigstartsignatures
\pysiglinewithargsret{\sphinxcode{\sphinxupquote{simvis.}}\sphinxbfcode{\sphinxupquote{visgen\_mwa\_multi}}}{\sphinxparam{\DUrole{n}{nside}}\sphinxparamcomma \sphinxparam{\DUrole{n}{grfmap}}\sphinxparamcomma \sphinxparam{\DUrole{n}{ra\_ptg}}\sphinxparamcomma \sphinxparam{\DUrole{n}{dec\_ptg}}\sphinxparamcomma \sphinxparam{\DUrole{n}{bl}}\sphinxparamcomma \sphinxparam{\DUrole{n}{nu}}}{}
\pysigstopsignatures
\sphinxAtStartPar
This calculates the visibilities for the baselines passed by bl array, given the RA, DEC, frequency, grfmap.
\begin{quote}\begin{description}
\sphinxlineitem{Parameters}\begin{itemize}
\item {} 
\sphinxAtStartPar
\sphinxstyleliteralstrong{\sphinxupquote{nside}} (\sphinxstyleliteralemphasis{\sphinxupquote{int}}) \textendash{} Healpix parameter. Usually power of 2.

\item {} 
\sphinxAtStartPar
\sphinxstyleliteralstrong{\sphinxupquote{grfmap}} (\sphinxstyleliteralemphasis{\sphinxupquote{np.ndarray}}) \textendash{} Contains the sky brightness temperature simulated from APS or PS for all the pixels.(Contains multiple realizations for 2D map or sky signal across all frequencies for 3D map). Shape \(N_{Realizations}\times 12\ {nside}^2\) for 2D, \(N_{c}\times 12\ {nside}^2\) for 3D.

\item {} 
\sphinxAtStartPar
\sphinxstyleliteralstrong{\sphinxupquote{ra\_ptg}} (\sphinxstyleliteralemphasis{\sphinxupquote{int}}\sphinxstyleliteralemphasis{\sphinxupquote{ or }}\sphinxstyleliteralemphasis{\sphinxupquote{float}}) \textendash{} RA of the pointing direction.(In degrees)

\item {} 
\sphinxAtStartPar
\sphinxstyleliteralstrong{\sphinxupquote{dec\_ptg}} (\sphinxstyleliteralemphasis{\sphinxupquote{int}}\sphinxstyleliteralemphasis{\sphinxupquote{ or }}\sphinxstyleliteralemphasis{\sphinxupquote{float}}) \textendash{} DEC of the pointing direction.(In degrees)

\item {} 
\sphinxAtStartPar
\sphinxstyleliteralstrong{\sphinxupquote{bl}} (\sphinxstyleliteralemphasis{\sphinxupquote{np.ndarray}}) \textendash{} The array conatins the components of the baselines defined by basis vectors \(\hat{\mathbf{e}}_1,\hat{\mathbf{e}}_2,\hat{\mathbf{e}}_3\).
Shape \(N_{Baselines}\times 3\).

\item {} 
\sphinxAtStartPar
\sphinxstyleliteralstrong{\sphinxupquote{nu}} (\sphinxstyleliteralemphasis{\sphinxupquote{float}}) \textendash{} Observing frequency \(\nu\) in \(\mathrm{MHz}\) unit.

\end{itemize}

\sphinxlineitem{Returns}
\sphinxAtStartPar
\sphinxstylestrong{vis} \textendash{} Contains the simulated visibility for all the baselines given by bl. (Along axis = 0 is the different realizations or frequencies, axis = 1 is the baselines). Shape \(N_{realizations}\ \times\ N_{Baselines}\) for 2D, \(N_{c}\times \ N_{Baselines}\) for 3D.

\sphinxlineitem{Return type}
\sphinxAtStartPar
np.ndarray

\end{description}\end{quote}

\clearpage

\end{fulllineitems}

\index{sim\_vis\_hdul() (in module simvis)@\spxentry{sim\_vis\_hdul()}\spxextra{in module simvis}}

\begin{fulllineitems}
\phantomsection\label{\detokenize{simvis:simvis.sim_vis_hdul}}
\pysigstartsignatures
\pysiglinewithargsret{\sphinxcode{\sphinxupquote{simvis.}}\sphinxbfcode{\sphinxupquote{sim\_vis\_hdul}}}{\sphinxparam{\DUrole{n}{hdulist}}\sphinxparamcomma \sphinxparam{\DUrole{n}{skysimtype}}\sphinxparamcomma \sphinxparam{\DUrole{n}{seed=None}}\sphinxparamcomma \sphinxparam{\DUrole{n}{apsfunc=\textless{}function aps\textgreater{}}}\sphinxparamcomma \sphinxparam{\DUrole{n}{psfunc=\textless{}function ps\textgreater{}}}}{}
\pysigstopsignatures
\sphinxAtStartPar
This calculates the visibilities for the baselines passed by bl, given the RA, DEC, frequency. It extracts the necessary information from the \sphinxcode{\sphinxupquote{hdulist}}, like the baseline distribution, centered frequency, frequency resolution, no of frequency channels, RA, DEC.
\begin{quote}\begin{description}
\sphinxlineitem{Parameters}\begin{itemize}
\item {} 
\sphinxAtStartPar
\sphinxstyleliteralstrong{\sphinxupquote{hdulist}} (\sphinxstyleliteralemphasis{\sphinxupquote{object}}) \textendash{} A fits object

\item {} 
\sphinxAtStartPar
\sphinxstyleliteralstrong{\sphinxupquote{skysimtype}} (\sphinxstyleliteralemphasis{\sphinxupquote{str}}) \textendash{} Either ‘2D’ or ‘3D’. If given ‘2D’ then creates multiple realizations. If given ‘3D’ then makes a 3D sky map.

\item {} 
\sphinxAtStartPar
\sphinxstyleliteralstrong{\sphinxupquote{seed}} (\sphinxstyleliteralemphasis{\sphinxupquote{int}}\sphinxstyleliteralemphasis{\sphinxupquote{ or }}\sphinxstyleliteralemphasis{\sphinxupquote{float}}) \textendash{} A seed value for generating GRF. (By deafult set to \sphinxcode{\sphinxupquote{None}})

\item {} 
\sphinxAtStartPar
\sphinxstyleliteralstrong{\sphinxupquote{apsfunc}} (\sphinxstyleliteralemphasis{\sphinxupquote{function}}) \textendash{} The input Angular Power Spectrum function. By default it is \(\texttt{aps}\) which has the form \(C_{\ell}=A\, \left(\frac{\ell}{\ell_0}\right)^{\beta}\) with \(A = 100\;\text{mK}^2\), \(\ell_0=1\), \(\beta\) = \sphinxhyphen{}2.

\item {} 
\sphinxAtStartPar
\sphinxstyleliteralstrong{\sphinxupquote{psfunc}} (\sphinxstyleliteralemphasis{\sphinxupquote{function}}) \textendash{} The input Power Spectrum function. By default it is \(\texttt{ps}\) which has the form \(A \left(\dfrac{k}{k_0}\right)^{s}\), with \(A=1, k_0=1, s=-2\).

\end{itemize}

\sphinxlineitem{Return type}
\sphinxAtStartPar
Simulates the visibilities for the input Power Spectrum or Angular Power Spectrum for the RA, DEC, frequency and baseline distribution given in the \sphinxcode{\sphinxupquote{hdulist}} object. (And puts the multiple realizations along the frequency axis if simulated sky signal is 2D).

\end{description}\end{quote}

\clearpage

\end{fulllineitems}

\index{sim\_vis() (in module simvis)@\spxentry{sim\_vis()}\spxextra{in module simvis}}

\begin{fulllineitems}
\phantomsection\label{\detokenize{simvis:simvis.sim_vis}}
\pysigstartsignatures
\pysiglinewithargsret{\sphinxcode{\sphinxupquote{simvis.}}\sphinxbfcode{\sphinxupquote{sim\_vis}}}{\sphinxparam{\DUrole{n}{uvinfits}}\sphinxparamcomma \sphinxparam{\DUrole{n}{uvoutfits}}\sphinxparamcomma \sphinxparam{\DUrole{n}{skysimtype}}\sphinxparamcomma \sphinxparam{\DUrole{n}{seed=None}}\sphinxparamcomma \sphinxparam{\DUrole{n}{apsfunc=\textless{}function aps\textgreater{}}}\sphinxparamcomma \sphinxparam{\DUrole{n}{psfunc=\textless{}function ps\textgreater{}}}}{}
\pysigstopsignatures
\sphinxAtStartPar
This calculates the visibilities for the baselines passed by bl\_file, given the RA, DEC, frequency. It extracts the necessary information from the \sphinxcode{\sphinxupquote{uvinfits}} file, like the baseline distribution, centered frequency, frequency resolution, no of frequency channels, RA, DEC and writes the simulated visibilities in a new fits file named \sphinxcode{\sphinxupquote{uvoutfits}} file.
\begin{quote}\begin{description}
\sphinxlineitem{Parameters}\begin{itemize}
\item {} 
\sphinxAtStartPar
\sphinxstyleliteralstrong{\sphinxupquote{uvinfits}} (\sphinxstyleliteralemphasis{\sphinxupquote{fits file}}) \textendash{} Input fits File.

\item {} 
\sphinxAtStartPar
\sphinxstyleliteralstrong{\sphinxupquote{uvoutfits}} (\sphinxstyleliteralemphasis{\sphinxupquote{fits file}}) \textendash{} Output fits File.

\item {} 
\sphinxAtStartPar
\sphinxstyleliteralstrong{\sphinxupquote{skysimtype}} (\sphinxstyleliteralemphasis{\sphinxupquote{str}}) \textendash{} Either ‘2D’ or ‘3D’. If given ‘2D’ then creates multiple realizations. If given ‘3D’ then makes a 3D sky map.

\item {} 
\sphinxAtStartPar
\sphinxstyleliteralstrong{\sphinxupquote{seed}} (\sphinxstyleliteralemphasis{\sphinxupquote{int}}\sphinxstyleliteralemphasis{\sphinxupquote{ or }}\sphinxstyleliteralemphasis{\sphinxupquote{float}}) \textendash{} A seed value for generating GRF.

\item {} 
\sphinxAtStartPar
\sphinxstyleliteralstrong{\sphinxupquote{apsfunc}} (\sphinxstyleliteralemphasis{\sphinxupquote{function}}) \textendash{} The input Angular Power Spectrum function. By default it is \(\texttt{aps}\) which has the form \(C_{\ell}=A\, \left(\frac{\ell}{\ell_0}\right)^{\beta}\) with \(A = 100\;\text{mK}^2\), \(\ell_0=1\), \(\beta\) = \sphinxhyphen{}2.

\item {} 
\sphinxAtStartPar
\sphinxstyleliteralstrong{\sphinxupquote{psfunc}} (\sphinxstyleliteralemphasis{\sphinxupquote{function}}) \textendash{} The input Power Spectrum function. By default it is \(\texttt{ps}\) which has the form \(A \left(\dfrac{k}{k_0}\right)^{s}\), with \(A=1, k_0=1, s=-2\).

\end{itemize}

\sphinxlineitem{Return type}
\sphinxAtStartPar
Simulates the visibilities for the input Power Spectrum or Angular Power Spectrum for the RA, DEC, frequency and baseline distribution given in the \sphinxcode{\sphinxupquote{uvinfits}} file. (And puts the multiple realizations along the frequency axis if simulated sky signal is 2D). Writes the data in \sphinxcode{\sphinxupquote{uvoutfits}} file.

\end{description}\end{quote}

\end{fulllineitems}




\renewcommand{\indexname}{Index}
\printindex
\end{document}